\chapter{参数敏感性与稳健性分析}

\section{分析目的}
在实际产线中，成本参数、故障识别能力和刀具寿命分布会随工况变化而波动。若最优策略只在单一参数组合下有效，则模型工程价值有限。因此本章围绕参数扰动开展系统敏感性分析，重点回答三个问题：
\begin{enumerate}
  \item 哪些参数对最优成本 $\phi^\star$ 影响最大；
  \item 最优检查节点是否具有结构稳定性；
  \item 在保守决策场景下，如何给出“稳健型”策略。
\end{enumerate}

\section{单因素敏感性}
设基准参数向量为
\[
\Theta_0=(\theta_1,\theta_2,\theta_3,\theta_4,\theta_5,R_1,R_2,\mu,\sigma).
\]
对任一参数 $\vartheta\in\Theta_0$，设扰动比例为 $\delta\in[-20\%,20\%]$，定义相对敏感度：
\[
S_{\vartheta}(\delta)=\frac{\phi^\star(\vartheta(1+\delta))-\phi^\star(\vartheta)}{\phi^\star(\vartheta)\delta}.
\]
若 $|S_{\vartheta}|$ 较大，则说明该参数对系统成本影响显著。

\subsection{废品损失参数 $\theta_1$}
当 $\theta_1$ 上升时，故障后继续生产的代价显著增加，最优策略会把检查节点整体前移，并倾向增加后段检查密度。该现象说明模型能够正确反映“质量风险上升 $\Rightarrow$ 预防加强”的管理直觉。

\subsection{检查成本参数 $\theta_2$}
当 $\theta_2$ 上升时，最优检查次数通常下降，但并不会无限减少。原因在于：若检查次数过低，废品成本会以更快速度抬升。二者相互制衡后形成新的最优点。

\subsection{恢复成本参数 $\theta_3$}
$\theta_3$ 上升会提高“故障后再恢复”的惩罚强度，进而提升预防性换刀的吸引力。在数值结果上表现为：最优终止检查点更早，主动换刀概率 $\eta_2$ 上升。

\section{双因素交互分析}
为观察交互效应，选取两组关键参数进行网格实验。

\subsection{$(\theta_1,\theta_2)$ 交互}
当 $\theta_1$ 高且 $\theta_2$ 低时，最优策略会显著加密检查；当 $\theta_1$ 低且 $\theta_2$ 高时，则会稀疏检查并依赖周期性换刀。说明“检查频度”受损失系数和检查成本共同控制，而非单参数单调决定。

\subsection{$(R_1,R_2)$ 交互}
在第二问设定下，$R_1$ 下降（正常品误判概率上升）与 $R_2$ 上升（故障品漏检概率上升）都会恶化决策质量。其交互结果通常表现为：
\begin{itemize}
  \item 单检策略成本曲线变平，最优解不再尖锐；
  \item 双检确认策略相对优势扩大；
  \item 最优检查次数可能从“中等频度”转向“前密后疏”的折中结构。
\end{itemize}

\section{分布参数扰动分析}
刀具寿命的均值与方差变化会改变故障位置分布。设
\[
\mu'=\mu(1+\delta_\mu),\quad \sigma'=\sigma(1+\delta_\sigma).
\]
有以下规律：
\begin{enumerate}
  \item $\mu$ 增大（寿命延长）时，整体检查节点后移；
  \item $\sigma$ 增大（离散性加大）时，后段检查密度增加；
  \item 当 $\sigma$ 过大时，固定间隔策略性能劣化明显，不等间隔策略优势增强。
\end{enumerate}

\section{稳健优化建议}
若企业希望减少策略频繁切换，可采用如下稳健决策准则：
\[
\min_{\pi\in\Pi}\left\{\mathbb E_{\Theta}[\phi(\pi,\Theta)]+\lambda\cdot \operatorname{Var}_{\Theta}[\phi(\pi,\Theta)]\right\},
\]
其中 $\pi$ 为策略，$\lambda$ 为风险厌恶系数。该准则兼顾平均成本与波动风险，适合多班组、多工况场景。

\section{本章小结}
敏感性实验表明：模型对关键参数变化具有可解释响应；最优策略在中等扰动范围内保持结构稳定；对于高不确定性场景，采用稳健目标函数可获得更稳定的管理方案。
